\chapter{Provision}
\label{chap:provision}

% Export diagrams/provision_flow.drawio → diagrams/provision_flow.pdf.
% Float drží celou stranu na diagram; text kapitoly má zůstat na jedné straně.
\begin{figure}[p]
  \centering
  \IfFileExists{../diagrams/provision_flow.pdf}{%
    \includegraphics[width=\textwidth,height=0.88\textheight,keepaspectratio]%
      {../diagrams/provision_flow.pdf}%
  }{%
    \fbox{%
      \parbox[c][0.82\textheight][c]{0.94\textwidth}{\centering}%
    }%
  }
  \caption{Rozhodovací strom příkazu \texttt{PROVISION}
    (\texttt{diagrams/provision\_flow}).}
  \label{fig:provision-flow}
\end{figure}

\noindent
Příkaz \texttt{PROVISION} naplní čip jednorázovými pady z QKD.
SE otevře TLS jen při \texttt{se\_\allowbreak ready\_\allowbreak provision()} (vlastník, certifikát
a klíč zařízení, SAE CA, ML-KEM klíč) a synchronizovaném čase. Po mTLS pošle
uplink v5 vázaný na TLS exporter (\texttt{SeSessionUplink.\allowbreak readAndVerify})
a souběžně smaže QKD sloty R-MEM; konec hlásí \texttt{TROPIC\_\allowbreak WIPE\_\allowbreak FINISHED}.
Obrázek~\ref{fig:provision-flow} je strom SaeNode: výběr peer SAE a peer
zařízení, lokální \texttt{synchronize()} a podle výsledku nová výměna
(\texttt{INSERTED}), uložený share (\texttt{ORIGIN\_PROVISIONED}, nebo
\texttt{RECORD\_SHARED\_WITH\_DIFFERENT\_SAE}), nebo konec bez klíčů
(\texttt{NOT\_INSERTED}).

Nová výměna v \texttt{PeerRecordSync} rezervuje pár
na obou SAE (\texttt{RECORD\_\allowbreak SYNCHRONIZATION}), zamkne fetch
(\texttt{RECORD\_\allowbreak FETCHING\_\allowbreak STARTED}) a stáhne pady z KME. Peer SAE je přes
RMI \texttt{insert} dohledá podle ID klíčů u svého KME a uloží zabalené v HSM.
Oba řádky pak přejdou do \texttt{RECORD\_\allowbreak SHARED}. Origin pošle downlink.
Po \texttt{SE\_\allowbreak OK} čipu zapíše \texttt{ORIGIN\_\allowbreak PROVISIONED} nejdřív na peeru,
pak u sebe. Až po návratu obou zápisů terminál vytiskne 100\,\% a na dalším
řádku \texttt{Share successful.}
Peer čip dostane materiál později, jen když jsou oba řádky
\texttt{ORIGIN\_\allowbreak PROVISIONED}. Peer SAE smaže nejdřív řádek na originu,
pak svůj, a teprve potom pošle downlink.

Mazání je vždy nejdřív vzdálené: rollback, \texttt{deleteShareAndRestart}
i doručení. Každé RMI volání jde přes mTLS s připnutým SPKI peeru. Volající
musí být vydávající SAE, nebo strana záznamu. Výjimkou je \texttt{deleteShare}
od jiného SAE (předání share), povolené jen pro dokončený share.

\section{Formální ověření}

PlusCal v \texttt{specs/sae-provision/SaeProvision.tla} sleduje tuto ságu
včetně pádu origin i peer JVM mezi každým krokem, vypršení pětiminutové
rezervace, výpadku spojení na origin nebo peer a jednoho pozdějšího pokusu
(nový fill, nebo doručení na peer).
TLS, PIN a bajty padů v modelu nejsou. Důkaz je prohledání stavů TLC podle
\texttt{SaeProvision.cfg}. Ve všech dosažitelných stavech platí:

\begin{itemize}
  \setlength{\itemsep}{0.15em}
  \item Oba committnuté share mají shodné příznaky materiálu.
  \item \texttt{ORIGIN\_PROVISIONED} na originu znamená totéž na peeru
        a oba vznikají jen po ACK origin čipu.
  \item Klíče na peer čipu znamenají klíče na origin čipu a už smazané řádky.
  \item Hláška úspěchu odpovídá konzistentnímu share, nebo jeho pozdějšímu
        spotřebování.
  \item Po čtení při dosažitelném originu i peeru za oknem rezervace je osiřelý
        \texttt{RECORD\_SHARED} pryč, nebo origin ještě drží čerstvou rezervaci.
\end{itemize}

\noindent
Stejné invarianty drží dva krajní běhy. Nedosažitelný origin nebo peer se
nemaže ani neplní, takže \texttt{RECORD\_SHARED} může zůstat. Pád po smazání řádků
a před zápisem na peer čip nechá klíče jen na originu, a stejný share se
nepošle dvakrát.
